% Section 9 — Threats to Validity (draft; final wording after results)
\section{Threats to Validity}
\label{sec:threats}

\subsection{Internal validity}
\label{sec:threats-internal}

\paragraph{Harness defects.} A broken harness can manufacture the result the author expects. We
therefore make the negative control (M0 under a crash point that must duplicate) and the positive
controls (M1-kafka and M4+M3 under fault-free runs, which must not) \emph{hard gates}: if either
control fails, the campaign halts and no measurement is analysed. The oracle is mutation-tested
before any experiment: synthetic defects (dropped row, duplicated row, flipped leg sign, altered
amount, spurious row) must be detected with 100\,\% sensitivity and attributed to the correct
invariant. The checker shares no code with the pipeline and recomputes ground truth from the
specification, so a bug in the service cannot propagate into its own scoring.

\paragraph{Faults that did not fire.} A fault-injection run is evidence only if the fault actually
occurred. Every injector appends a timestamped receipt to \path{results/faults.log} when it fires
and when it is restored; runs without a complete receipt pair are excluded and counted, and the
exclusion rate is reported per cell. This converts ``the fault fired rarely'' from a hidden threat
into a measured quantity.

\paragraph{Configuration bias.} Mechanisms differ in their natural configuration choices, and a
comparison that tunes one mechanism more carefully than another is unfair. We pin every
non-intrinsic parameter identically (consumer, producer, JDBC, JVM, broker, database settings),
allow only mechanism-intrinsic differences (deduplication index, outbox tables, progress table,
transactional producer), hash the effective configuration per run, and fail the campaign if the
shared hash differs between mechanisms. Mechanism implementations live in one code base and share
the poll/commit skeleton.

\paragraph{Order drift.} Long campaigns drift: thermal state, page cache, background tasks and
disk fragmentation all change over hours. Runs are randomised and interleaved across mechanisms
rather than blocked by mechanism, paired seeds are used within a cell, the full stack is restarted
at a fixed cadence, and run order is recorded so that drift can be tested as a covariate.

\subsection{Construct validity}
\label{sec:threats-construct}

\paragraph{A simplified ledger.} We model one currency, balanced two-leg entries, append-only
postings and no reversals or claim adjustments. Real ledgers have multi-leg entries across
currencies, corrections and reconciliation states. These add failure surface but do not change the
atomicity question we study: an external posting either matches a consumed event exactly once or
it does not. The oracle deliberately tests the simple invariant set (I1--I4) because these are the
properties a ledger must never violate, and they can be checked exactly.

\paragraph{The oracle versus accounting practice.} Our oracle measures the invariants we defined,
not the full correctness of an accounting system. A pipeline can pass our checks while still
being wrong about business rules (e.g.\ a mis-classified account). We therefore report the
invariant-level result and do not claim end-to-end accounting correctness. Conversely, the oracle
can only detect what it is pointed at: an effect not represented in the expected posting set
cannot be scored. The mutation-testing gate is the mitigation and is reported with the experiment.

\paragraph{Synthetic load.} Events are generated deterministically from a seed with a fixed size
distribution and a balanced entry structure. Real traffic has hot keys, large entries, bursts and
cross-record dependencies. Our design intentionally removes these so that correctness outcomes are
attributable to the mechanism rather than to workload variance; the performance numbers are
therefore lower bounds on the overhead differences, not absolute production capacities.

\subsection{External validity}
\label{sec:threats-external}

\paragraph{Single host, native processes.} The reported experiments run on one Windows host as
native processes: three Kafka KRaft brokers, one PostgreSQL cluster and Toxiproxy. There is no
container isolation, no multi-AZ placement and no real machine failure; a broker crash is a process
kill. Process kills and TCP cuts are real faults at the OS level, but they do not reproduce rack
power loss, fsync loss on a virtualised disk, or coordinated failures. The direction of the bias
is toward \emph{understating} the failure surface: a single host has no network partitions between
Kafka replicas, so replicated durability assumptions are never tested by the environment itself.
All conclusions are therefore about the modelled mechanism semantics, whose failure windows are
explicitly enumerated, and not about a specific cloud deployment.

\paragraph{Version specificity.} Results are tied to the pinned versions in
\code{tools/TOOLCHAIN-VERSIONS.txt} (Kafka 3.7.2, PostgreSQL 16.11, JDK 21). Kafka's transaction
protocol has been stable since KIP-98/KIP-447, but default settings (rebalance protocol, producer
batching) change between minor releases. The artifact pins versions and records digests so the
experiment can be repeated exactly; claims are stated at the semantic level so that they survive
version upgrades, while measured overheads are reported with their configuration.

\subsection{Conclusion validity}
\label{sec:threats-conclusion}

Correctness outcomes are mostly-zero counts, and ``zero observed'' is never ``proved''. We report
the proportion of violating runs with Clopper--Pearson 95\,\% confidence intervals and, for
zero-event cells, the rule-of-three upper bound, so the reader can see how much evidence a zero
carries. Performance comparisons use pre-registered non-parametric tests (Friedman with Nemenyi
post-hoc), paired tests where the design is paired, Holm correction within each family, and
Cliff's delta with confidence intervals; analyses not in the pre-registration are labelled
exploratory. Multiple endpoints (correctness, throughput, latency, recovery time) could inflate
false positives; we fix the families and corrections before looking at the data. The effective
sample size is bounded by the wall-clock budget; the staged design and the cut ladder are recorded
in \code{results/DEVIATIONS.md} so that the reader knows exactly which cells were executed and
which were reduced.

\subsection{Reproducibility}
\label{sec:threats-repro}

Every quantitative claim traces to a run record produced by a logged command with UTC timestamp,
git commit, host and seeds (\code{results/RUNLOG.md}), and every table and figure is regenerated
by scripts from committed raw data. Ground-truth files are regenerated deterministically from the
master seed and verified against the pre-send digest; configurations are hashed; analyses are
frozen before the confirmatory stage. The container-based reproduction path is provided for CI,
and the native-process path used for the reported runs is scripted and version-pinned. What cannot
be reproduced bit-for-bit is timing on a different host: we therefore report distributions and
paired comparisons rather than absolute single numbers, and the threats section states the host
configuration.
